% ==============================================================================
% capitulos/02_fundamentacao.tex
% Capítulo 2 — Fundamentação Teórica e Tecnologias Adotadas
% ==============================================================================

\chapter{Fundamentação Teórica e Tecnologias Adotadas}
\label{cap:fundamentacao}

Este capítulo estabelece a base conceitual e tecnológica que sustenta o projeto, abrangendo o controlador industrial adotado, o protocolo de comunicação de rede, a plataforma Node-RED de programação baseada em fluxos, os fundamentos da camada analítica híbrida e os requisitos de separação entre segurança operacional e supervisão inteligente.

\section{Controladores Lógicos Programáveis e o WEG PLC500}
\label{sec:fund_plc500}

Os Controladores Lógicos Programáveis (CLPs) são sistemas computacionais industriais projetados para operar em ciclos determinísticos de varredura (\textit{scan cycle}), executando sequencialmente a leitura das entradas físicas, o processamento da lógica de controle, a atualização das saídas físicas e as rotinas de diagnóstico e comunicação \cite{georgini_automacao, moraes_castrucci}. A norma internacional \textbf{IEC 61131-3} padroniza a estrutura de software e as linguagens de programação desses equipamentos, contemplando linguagens gráficas --- como \textit{Ladder Diagram} (LD) e \textit{Function Block Diagram} (FBD) --- e linguagens textuais, entre as quais se destaca o \textit{Structured Text} (ST) pela expressividade na construção de máquinas de estados, temporizações, manipulação de registradores e rotinas de intertravamento \cite{iec_61131_3}.

O WEG PLC500 é um controlador de médio e alto desempenho baseado no ambiente CODESYS V3.5, projetado para controle de máquinas, automação de processos e conectividade IIoT \cite{weg_plc500_manual, codesys_v35}. Seu hardware dispõe de entradas e saídas digitais nativas suficientes para a implementação integral da bancada didático-industrial, além de suportar expansão modular por cartões da série MOD (compartilhada com as unidades remotas RUW100). O Quadro~\ref{quadro:spec_plc500} resume suas especificações técnicas principais.

\begin{quadro}[htb]
  \centering
  \SingleSpacing
  \caption{Especificações técnicas principais do controlador WEG PLC500 aplicadas ao projeto}
  \label{quadro:spec_plc500}
  \footnotesize
  \begin{tabularx}{\textwidth}{|p{4.6cm}|L|}
    \hline
    \textbf{Característica} & \textbf{Especificação Técnica (WEG PLC500)} \\
    \hline
    Arquitetura de processamento & Processador i.MX Dual Core de 1\,GHz associado a coprocessador de 240\,MHz \\
    \hline
    Memória embarcada & 1\,GB de memória RAM e 4\,GB de memória Flash não volátil \\
    \hline
    Entradas digitais nativas (DI) & 8 entradas digitais PNP (24\,Vcc), sendo 4 entradas rápidas de até 150\,kHz \\
    \hline
    Saídas digitais nativas (DO) & 8 saídas digitais PNP (24\,Vcc) para acionamento de cargas discretas \\
    \hline
    Interfaces Ethernet & 2 portas RJ45 Gigabit Ethernet independentes (10/100/1000\,Mbps: \texttt{ETH1} e \texttt{ETH2}) \\
    \hline
    Expansibilidade local & Até 8 módulos de expansão \textit{Plug and Play} da família MOD (compatíveis com RUW100) \\
    \hline
    Ambiente de programação & CODESYS V3.5 (padrão IEC 61131-3: ST, LD, FBD, SFC e CFC) \\
    \hline
    Suporte Modbus TCP/IP & Nativo no CODESYS como \textit{Modbus TCP Server} (Slave) e \textit{Modbus TCP Client} (Master) \\
    \hline
  \end{tabularx}
  \fonte{Adaptado de \citeonline{weg_plc500_manual} e \citeonline{weg_plc500_modbus}.}
\end{quadro}

Um diferencial arquitetural do WEG PLC500 para projetos de convergência OT/IT é a disponibilidade de duas interfaces Ethernet físicas independentes (\texttt{ETH1} e \texttt{ETH2}), permitindo segregar a rede de engenharia CODESYS da sub-rede dedicada à comunicação Modbus TCP/IP com o Node-RED \cite{weg_plc500_modbus, stouffer_nist_ot}.

\section{Protocolo Modbus TCP/IP}
\label{sec:fund_modbus}

Desenvolvido originalmente para barramentos seriais, o protocolo Modbus tornou-se o padrão aberto mais difundido em automação industrial devido à sua simplicidade e baixa sobrecarga de processamento. Em redes Ethernet, sua implementação segue duas especificações da Modbus Organization: o \textit{Modbus Application Protocol Specification V1.1b3} \cite{modbus_app_v11b3}, referente à camada de aplicação, e o \textit{Modbus Messaging on TCP/IP Implementation Guide V1.0b} \cite{modbus_tcp_v10b}, que define o encapsulamento cliente/servidor sobre conexões TCP na porta padrão 502.

No Modbus TCP/IP, a Unidade de Dados de Aplicação (ADU) substitui o endereço serial e o verificador CRC por um cabeçalho MBAP (\textit{Modbus Application Protocol}) de 7 bytes anteposto à Unidade de Dados de Protocolo (PDU):
\begin{itemize}
  \item Transaction Identifier (2 bytes): identifica e pareia requisições e respostas sobre o mesmo \textit{socket} TCP;
  \item Protocol Identifier (2 bytes): preenchido com \texttt{0x0000} para transações Modbus padrão;
  \item Length (2 bytes): informa o número de bytes subsequentes do pacote (\textit{Unit Identifier} somado à PDU);
  \item Unit Identifier (1 byte): identifica o dispositivo lógico de destino (configurado como \texttt{1} para acesso direto ao CLP).
\end{itemize}

A memória exposta pelo servidor Modbus divide-se em quatro tabelas primárias acessadas por Códigos de Função (\textit{Function Codes} -- FC), conforme sistematizado no Quadro~\ref{quadro:modbus_modelo}. Requisições inválidas recebem como resposta um quadro de exceção com o código de função somado a \texttt{0x80} e o respectivo código diagnóstico de erro \cite{modbus_app_v11b3}.

\begin{quadro}[htb]
  \centering
  \SingleSpacing
  \caption{Modelo de dados primário e códigos de função Modbus TCP/IP adotados no projeto}
  \label{quadro:modbus_modelo}
  \footnotesize
  \begin{tabularx}{\textwidth}{|p{2.9cm}|p{1.6cm}|p{1.7cm}|p{2.3cm}|L|}
    \hline
    \textbf{Tabela Modbus} & \textbf{Tipo} & \textbf{Acesso} & \textbf{Códigos (FC)} & \textbf{Papel na Arquitetura} \\
    \hline
    \textit{Discrete Inputs} & 1 bit & Leitura & \texttt{FC02} & Estados filtrados das botoeiras (\texttt{DI0}--\texttt{DI3}) e flags de segurança do PLC500 \\
    \hline
    \textit{Coils} & 1 bit & Leitura / Escrita & \texttt{FC01}, \texttt{FC05}, \texttt{FC15} & Comandos supervisórios do Node-RED para atuação nos LEDs e \textit{reset} de falha \\
    \hline
    \textit{Input Registers} & 16 bits & Leitura & \texttt{FC04} & Contadores acumulados de acionamentos, estado da máquina do CLP e telemetria real \\
    \hline
    \textit{Holding Registers} & 16 bits & Leitura / Escrita & \texttt{FC03}, \texttt{FC06}, \texttt{FC16} & \textit{Heartbeat} do Node-RED, \textit{timeout} do \textit{watchdog} e severidade classificada pela IA \\
    \hline
  \end{tabularx}
  \fonte{Adaptado de \citeonline{modbus_app_v11b3} e \citeonline{modbus_tcp_v10b}.}
\end{quadro}

\section{Programação Baseada em Fluxos com Node-RED e Integração Modbus}
\label{sec:fund_nodered}

O Node-RED, mantido pela OpenJS Foundation, é um ambiente de programação visual orientado a eventos construído sobre o motor assíncrono do Node.js, amplamente adotado em gateways IIoT e sistemas supervisórios leves \cite{nodered_docs}. Dentro de um fluxo (\textit{flow}), o processamento ocorre pela passagem de objetos JavaScript estruturados (mensagens \texttt{msg}, cujo conteúdo principal reside em \texttt{msg.payload}) através de nós especializados, contando ainda com áreas de memória persistente em escopo de nó (\texttt{context}), de aba (\texttt{flow}) ou global (\texttt{global}).

A integração nativa do Node-RED com redes industriais Modbus TCP/IP e Serial é realizada pelo pacote \texttt{node-red-contrib-modbus}, desenvolvido e mantido por Klaus Landsdorf \cite{nodered_contrib_modbus}. Esse pacote fornece um nó de configuração compartilhado (\texttt{Modbus-Client}) --- responsável por manter a conexão TCP persistente na porta~502, gerenciar reconexões automáticas e enfileirar transações sequencialmente --- acompanhado de nós operacionais dedicados:
\begin{itemize}
  \item \texttt{Modbus-Read} e \texttt{Modbus-Flex-Getter}: executam leituras cíclicas ou disparadas por mensagem (\texttt{FC01} a \texttt{FC04}), entregando na saída tanto o vetor de valores decodificados em \texttt{msg.payload} quanto o buffer bruto em \texttt{msg.responseBuffer};
  \item \texttt{Modbus-Write} e \texttt{Modbus-Flex-Write}: permitem escrever bits individuais ou múltiplos (\texttt{FC05}, \texttt{FC15}) e registradores de 16~bits (\texttt{FC06}, \texttt{FC16}) dinamicamente a partir dos parâmetros montados por um nó \texttt{function} anterior (\texttt{msg.payload.value}, \texttt{msg.payload.fc}, \texttt{msg.payload.address} e \texttt{msg.payload.quantity}).
\end{itemize}

Associado a nós \texttt{function} (JavaScript), conectores para banco de dados PostgreSQL (\texttt{node-red-contrib-postgresql}) e nós \texttt{http request} para consumo de APIs de modelos de linguagem, o Node-RED reúne em uma única ferramenta determinismo de coleta e flexibilidade analítica \cite{nodered_docs, postgresql_docs}.

\section{Análise Híbrida de Dados: Estatística Temporal, Aprendizado de Máquina e Agentes LLM}
\label{sec:fund_ia}

A detecção de anomalias em processos industriais discretos requer identificar desvios em relação ao padrão operacional esperado, como repiques ou oscilações anômalas de sensores (\textit{chattering}), ativação prolongada de chaves de nível ou ocorrências térmicas críticas \cite{chandola_2009}. Submeter amostras brutas coletadas a cada $500\,\text{ms}$ diretamente a um LLM resultaria em latência elevada, alto custo por requisição e baixa reprodutibilidade matemática.

Para contornar essa limitação, implementa-se uma arquitetura analítica híbrida dividida em três estágios cooperativos entre o Node-RED e um microsserviço Python (FastAPI / scikit-learn):
\begin{enumerate}
  \item Estágio estatístico determinístico (nó \texttt{function}): executado sobre janelas temporais deslizantes de duração $W = 60\,\text{s}$ com passo (\textit{stride}) de $15\,\text{s}$, extrai um vetor canônico de $10$ atributos (\texttt{MODEL\_FEATURE\_COLUMNS}) --- incluindo taxa de acionamento $\lambda_i(W)$, fração de tempo ativo (\textit{duty cycle}) $D_i(W)$, escore padronizado \textit{Z-score} $z_i(W)$, índice de trepidação (\texttt{chattering\_index}) e indicadores locais de veto --- utilizando a linha de base histórica armazenada no PostgreSQL (\texttt{vw\_baseline\_24h}) e em memória de contexto (\texttt{flow}).
  \item Estágio de aprendizado de máquina em tempo real (microsserviço FastAPI / scikit-learn): recebe o vetor de $10$ atributos da janela (\texttt{POST /predict}) ou o lote de amostras brutas (\texttt{POST /predict/raw-window}) e executa simultaneamente uma pipeline não supervisionada (\texttt{StandardScaler + IsolationForest}, com ROC-AUC de teste de $0{,}8993$) e um classificador supervisionado multiclasse (\texttt{HistGradientBoostingClassifier}, selecionado frente ao \texttt{RandomForestClassifier} com acurácia de teste de $99{,}73\%$ e Macro-F1 de $0{,}9956$). Os modelos são treinados sobre um conjunto sintético de $18\,\text{h}$ a $2\,\text{Hz}$ ($129.600$ amostras brutas e $3.688$ janelas deslizantes de $60\,\text{s}$) particionado cronologicamente em $65\%/15\%/20\%$ (treino/validação/teste) com lacunas de purga (\textit{purge gaps}) $\ge 60\,\text{s}$ para eliminar vazamento temporal entre partições.
  \item Estágio cognitivo baseado em agente LLM (nó \texttt{http request} + validador JSON): acionado seletivamente quando os estágios anteriores detectam anomalia, mudança de estado ou fechamento de turno. O agente analisa o vetor de métricas da janela, a predição do serviço ML e o histórico recente, gerando diagnóstico em linguagem natural, hipótese de causa provável, recomendação de comando para os LEDs e relatório técnico automático.
\end{enumerate}

Toda resposta proveniente do modelo passa obrigatoriamente por validação contra o JSON Schema Draft 2020-12 (via motor \texttt{Ajv}) e por um guardião determinístico de monotonicidade de segurança no Node-RED antes de alcançar o nó \texttt{Modbus-Flex-Write}; qualquer violação estrutural, tentativa de rebaixamento de severidade ou \textit{timeout} descarta a saída cognitiva e aplica imediatamente o comando determinístico seguro de \textit{fallback}.

\section{Segurança Funcional e Separação de Responsabilidades OT/IT}
\label{sec:fund_seguranca}

Em conformidade com as boas práticas de segurança de sistemas de controle industrial estabelecidas no guia NIST SP 800-82 Rev.\,3 \cite{stouffer_nist_ot} e na norma IEC 62443 \cite{iec_62443}, a proteção física de equipamentos e operadores jamais deve depender exclusivamente de fluxos executados na camada IT ou de enlaces Modbus TCP/IP sem autenticação criptográfica nativa.

Assim, a arquitetura fundamenta-se em três barreiras de defesa em profundidade:
\begin{itemize}
  \item Precedência local absoluta no WEG PLC500: eventos críticos de campo --- abertura do contato normalmente fechado (NF) da botoeira de parada de emergência ou ativação do sensor de sobretemperatura --- forçam imediatamente o estado seguro nas saídas físicas do CLP, vetando qualquer comando divergente enviado pelo Node-RED.
  \item Supervisão contínua de enlace por \textit{Watchdog} de \textit{Heartbeat}: o fluxo do Node-RED incrementa ciclicamente um registrador de pulso de vida no PLC500; se o valor deixar de ser atualizado por intervalo superior a $T_{wd}$, o CLP assume perda de comunicação, desliga o atuador principal e aciona sinalização local de falha.
  \item Segmentação de rede por interfaces físicas distintas: a comunicação Modbus TCP/IP opera isolada na interface \texttt{ETH2} do PLC500, com regras de firewall restringindo conexões na porta TCP~502 exclusivamente ao endereço IP do servidor Node-RED \cite{stouffer_nist_ot}.
\end{itemize}
